\section{An actual infinite-refresh counterexample}\label{app:trap}
This section proves the main counterexample for the original uncapped $K$ protocol. All logarithms are natural and terminating decimals in bounds are exact rational numbers. The finite multiplicities are specified analytically; no enormous feature matrix or trajectory is claimed to have been numerically executed. Labels (T1)--(T56) refer only to this construction. The proof follows every actual acceptance, variable step and refresh.
\subsection{The feature matrix}

Set

\[
R=100,\qquad B=\frac{101}{100},\qquad \delta=\frac{99}{100},
\]

and

\[
N=2^{10^7},\qquad m=N^3+N^2+N+1.
\]

Take nuisance dimension \(p=1\) and the following rows:

\begin{center}
\begin{tabular}{lll}\toprule
Type & Feature row & Multiplicity \\ \midrule
Base & \((1,0)\) & \(N^3\) \\
Auxiliary & \((\sqrt{\delta},0)\) & \(N^2\) \\
First informative row & \((\sqrt B,-R)\) & \(1\) \\
Second informative row & \((\sqrt B,1)\) & \(N\) \\
\bottomrule\end{tabular}
\end{center}

Initialize every individual row with weight \(1/m\), and cache exactly there.

The feature matrix has rank \(2\). Write the four \textbf{aggregate} weights as

\[
(w_0,\omega,x_1,x_2).
\]

Their initial values are

\[
(w_0,\omega,x_1,x_2)
=\frac1m(N^3,N^2,1,N).
\tag{T1}
\]

\subsubsection{Aggregation preserves the exact protocol}

Copies of the same row always have the same weight. If a row type has \(n_j\) copies and aggregate current/cache weights \(W_j,S_j\), then each copy has weights \(W_j/n_j,S_j/n_j\), so

\[
\frac{w_i}{s_i}=\frac{W_j}{S_j},
\qquad
\frac{d_i}{w_i}=\frac{p_j^2}{q}-1.
\]

All sums defining \(A_s,q,r,h,Q,H,g\) aggregate exactly, and so do the extrema defining \(\alpha,\gamma,\zeta\). Therefore the four-weight analysis below is exactly the specified algorithm on the \(m\)-row matrix.

For positive aggregate weights,

\[
A_w=R^2x_1+x_2>0,
\]

and direct elimination gives

\[
v(w)
=w_0+\delta\omega+
B(R+1)^2\frac{x_1x_2}{R^2x_1+x_2}.
\tag{T2}
\]

The optimum is

\[
\boxed{v^*=B.}
\tag{T3}
\]

Indeed, evaluation at \(\theta=0\) gives \(v(w)\le B\) for every design. Conversely, the design supported on the informative pair with

\[
x_1=\frac1{R+1},\qquad x_2=\frac R{R+1}
\]

has minimizing coefficient \(\theta=0\) and value \(B\).

Also, the stopping test never fires: the two positive-weight, zero-nuisance row types always have distinct cached squared residuals \(1\) and \(\delta\). Thus \(g>0\) at every iterate, including every refresh.



\subsection{An epoch lemma for the actual rule}

An epoch starts with a fresh cache and ends at the first state where that old cache fails \(G\ge g/2\), before refreshing there.

During an epoch, label the two informative rows as follows:

\(x\) is the weight of the row that is initially much smaller;\par
\(y\) is the weight of the other informative row.\par

The labels remain fixed within the epoch. There are two types of epoch: \(\mathrm H\), when \(x=x_1\), and \(\mathrm L\), when \(x=x_2\).

Define

\[
\mathcal A_{\mathrm H}=B(R+1)^2=\frac{1030301}{100},
\qquad
\mathcal A_{\mathrm L}=B(1+1/R)^2=\frac{1030301}{10^6},
\tag{T4}
\]

and

\[
a_\sigma=\log\frac{1+\mathcal A_\sigma}{1+\delta},
\qquad
\beta=\log(1+\delta),
\qquad
C_\sigma=\frac{a_\sigma+2\beta}{2a_\sigma+\beta}.
\tag{T5}
\]

\subsubsection{Epoch lemma}

Suppose a fresh epoch starts with

\[
\omega_s\le10^{-6},\qquad
-\log\frac{s_y}{\omega_s}=b\ge1000,
\qquad
-\log\frac{s_x}{\omega_s}=2b+\varepsilon,
\quad |\varepsilon|\le100.
\tag{T6}
\]

Then this epoch ends after finitely many steps. Throughout the epoch, including its ending state,

\[
\omega\le10^{-6},
\qquad
\frac{x}{\omega},\frac{y}{\omega}\le10^{-20}.
\tag{T7}
\]

The auxiliary weight \(\omega\) decreases, the first step is exactly \(1/2\), and every accepted step is at least \(10^{-12}\).

At the ending state, put

\[
b'=-\log\frac{x_{\rm end}}{\omega_{\rm end}}.
\]

Then

\[
\left|
-\log\frac{y_{\rm end}}{\omega_{\rm end}}-2b'
\right|\le100,
\qquad
|b'-C_\sigma b|\le1000.
\tag{T8}
\]

Consequently, whenever \(b'\ge1000\), interchanging the labels of the informative rows gives the hypotheses for the next epoch.

Here is a proof with explicit estimates.

\subsubsection{2.1. Cached residuals}

Write

\[
\rho=\frac{s_x}{s_y}=e^{-b-\varepsilon},
\qquad
\kappa=\frac{z_x^2}{z_y^2}\in\{10^4,10^{-4}\}.
\]

The two informative rows have opposite nuisance signs. Solving the scalar normal equation gives their cached squared residuals

\[
P:=p_x^2=\frac{\mathcal A_\sigma}{(1+\kappa\rho)^2},
\qquad
S:=p_y^2=\kappa\rho^2P.
\tag{T9}
\]

Moreover,

\[
A_s=z_y^2s_y(1+\kappa\rho).
\tag{T10}
\]

Since \(b\ge1000\) and \(|\varepsilon|\le100\),

\[
1.02<P<11000,\qquad 0<S<10^{-20},
\qquad 1\le1+\kappa\rho\le2.
\tag{T11}
\]

Thus the cached squared sensitivities are strictly ordered:

\[
S<\delta<1<P.
\tag{T12}
\]

\subsubsection{2.2. Estimates in the region (T7)}

First derive estimates conditional on (T7); an induction below establishes that the trajectory remains there.

Because \(w_0=1-\omega-x-y\),

\[
q=1-(1-\delta)\omega+(P-1)x+(S-1)y.
\]

Consequently,

\[
1-.011\omega\le q\le1-.009\omega,
\qquad .999<q<1.
\tag{T13}
\]

The identity

\[
g=\frac1q\sum_i w_i(p_i^2-q)^2
\]

then gives

\[
10^{-5}\omega\le g\le10^{-3}\omega.
\tag{T14}
\]

For the lower bound, the auxiliary row alone contributes
\(\omega(q-\delta)^2/q\), and \(q-\delta>.009\).
For the upper bound, the auxiliary contribution is below
\(0.000101\omega\), the base-row contribution is at most
\(.011^2\omega^2/.999\), and the two informative contributions are bounded using \(x,y\le10^{-20}\omega\) and \(P<11000\).

For a step \(t\), write

\[
F_U=1+t(U/q-1)
\]

for the multiplier of a row with cached squared residual \(U\). For \(0<t\le1/2\), (T11)–(T13) imply

\[
\frac{F_P}{F_\delta}<10^4,\qquad
\frac{F_S}{F_\delta}<1,\qquad
\frac{F_P^2}{F_\delta F_S}<10^{10},
\tag{T15}
\]

and

\[
F_\delta\le1-.009t.
\tag{T16}
\]

For example, \(F_P<5600\), \(F_\delta\ge.995\), and \(F_S\ge1/2\).
When \(t=1/2\), (T16) gives \(F_\delta<.996\).

Introduce the positive balance variable

\[
W=\frac{x^2}{\omega y}.
\tag{T17}
\]

\subsubsection{2.3. The first step really is \texorpdfstring{$1/2$}{one half}}

At the fresh state,

\[
W_0=\frac{s_x^2}{\omega_s s_y}
=e^{-3b-2\varepsilon}\le e^{-2800}.
\tag{T18}
\]

The cache normal equation gives

\[
z_xp_xs_x=-z_yp_ys_y.
\]

Hence, at the fresh state,

\[
H=
\frac{\kappa P}{1+\kappa\rho}
\frac{s_x^2}{s_y}
\left(\frac{P-S}{q}\right)^2.
\]

Using (T11), (T13), and (T14),

\[
\frac Hg\le 2\cdot10^{22}W_0<\frac12.
\tag{T19}
\]

Since fresh \(Q=0\), \(G=g\), and \(K=H\), the prescribed step is exactly

\[
t=\frac12.
\tag{T20}
\]

After this step, the relative current-to-cache growth of the two informative rows satisfies

\[
J:=\frac{x/s_x}{y/s_y}
=\frac{q+P}{q+S}>2.
\tag{T21}
\]

It increases strictly at all subsequent steps in this epoch.

\subsubsection{2.4. Exact control of \texorpdfstring{$\alpha,\gamma,\zeta,G,K$}{alpha, gamma, zeta, G, K}}

Let

\[
\eta_x=x/s_x,\qquad \eta_y=y/s_y.
\]

Because of the fixed ordering (T12), the products of update multipliers preserve the same ordering of current-to-cache ratios. In particular,

\[
\boxed{\alpha=\eta_y.}
\tag{T22}
\]

This includes the base and auxiliary rows and all their copies.

Let

\[
\mathcal C=z_xp_xs_x=-z_yp_ys_y,
\qquad
\Delta_x=P/q-1,\quad \Delta_y=S/q-1.
\]

At every nonfresh state,

\[
r=\mathcal C(\eta_x-\eta_y),
\qquad
h=\mathcal C(\eta_x\Delta_x-\eta_y\Delta_y).
\]

Therefore, using \(J>2\),

\[
\ell:=\frac hr
=\Delta_x+\frac{\Delta_x-\Delta_y}{J-1}
\quad\text{satisfies}\quad
0<\ell<10^5.
\tag{T23}
\]

Thus \(\sqrt H=\ell\sqrt Q\). Also,

\[
\gamma<10^5,\qquad
\zeta=1-S/q\in[1/2,1].
\tag{T24}
\]

A direct substitution into \(Q=r^2/A_s\), using (T10) and (T22), gives the exact formula

\[
\frac Q\alpha
=
\frac{\kappa P}{1+\kappa\rho}
(1-J^{-1})^2\frac{x^2}{y}.
\tag{T25}
\]

The coefficient multiplying \(x^2/y\) lies between \(10^{-6}\) and \(10^9\). Combining this with (T14), and defining

\[
\mathcal Z=\frac{Q}{\alpha g},
\]

yields

\[
10^{-3}W\le\mathcal Z\le10^{14}W.
\tag{T26}
\]

Crucially, the actual certificate quantities now satisfy

\[
1-\frac Gg=(2\ell+\zeta)\mathcal Z,
\qquad
\frac Kg=(\ell+\gamma)^2\mathcal Z.
\tag{T27}
\]

In particular,

\[
1-\frac Gg\le10^6\mathcal Z,
\qquad
\frac Kg\le4\cdot10^{10}\mathcal Z.
\tag{T28}
\]

At any accepted nonfresh state, \(G/g\ge1/2\). Since
\(2\ell+\zeta\ge1/2\), (T27) implies \(\mathcal Z\le1\). Therefore

\[
W\le10^3
\tag{T29}
\]

and the prescribed step satisfies

\[
t=\min\left(\frac12,\frac{G}{4K}\right)
\ge
\min\left(\frac12,\frac1{32\cdot10^{10}}\right)
>10^{-12}.
\tag{T30}
\]

No lower bound on any weight was used.

\subsubsection{2.5. Invariance of the region and finiteness of the epoch}

The starting state satisfies (T7). Its fresh half-step remains in that region by (T15), (T18).

For the subsequent induction, \(y/\omega\) decreases by (T15), so

\[
y/\omega\le e^{-b}.
\]

At every accepted nonfresh state, (T29) gives

\[
\frac{x}{\omega}
=\sqrt{W\,\frac y\omega}
\le\sqrt{10^3}\,e^{-b/2}.
\]

After the next step, (T15) therefore gives

\[
\frac{x_{\rm next}}{\omega_{\rm next}}
\le10^4\sqrt{10^3}\,e^{-b/2}
<10^{-20}.
\]

The other ratio also remains below \(10^{-20}\), and \(\omega\) decreases by (T16). This proves (T7) inductively, including the first failed state.

To prove that failure occurs after finitely many steps, define

\[
a_k=\log(F_P/F_\delta)>0,
\qquad
\beta_k=\log(F_\delta/F_S)>0.
\]

Then

\[
\log W_{k+1}-\log W_k=2a_k+\beta_k.
\tag{T31}
\]

Moreover,

\[
F_\delta-F_S=\frac{t(\delta-S)}q>.98t,\qquad F_S\le1,
\]

so

\[
\beta_k\ge\log(1+.98t)\ge\frac t2.
\tag{T32}
\]

Every accepted step thus increases \(\log W\) by at least
\(5\cdot10^{-13}\), whereas every accepted state has \(W\le10^3\), including the fresh state. Infinite retention of this cache is impossible.

For completeness, an explicit upper bound on the number \(M\) of steps in the epoch is

\[
M\le
2+\left\lceil
2\cdot10^{12}\bigl(3b+200+\log(10^3)\bigr)
\right\rceil.
\tag{T33}
\]

Thus no unknown hitting time is being assumed.

At the ending state, \(g>0\) and \(G/g<1/2\). Equations (T26) and (T28) give

\[
\frac12<10^6\mathcal Z\le10^{20}W_{\rm end},
\]

so \(W_{\rm end}>10^{-22}\). The preceding accepted state had \(W\le10^3\), and the one-step multiplier in (T15) is below \(10^{10}\). Hence

\[
10^{-22}<W_{\rm end}\le10^{13}.
\tag{T34}
\]

In particular,

\[
|\log W_{\rm end}|\le100.
\tag{T35}
\]

\subsubsection{2.6. The logarithmic epoch map, with a uniform error bound}

It remains to prove the second estimate in (T8). This requires controlling the actual variable steps, not replacing them by half-steps.

Whenever a nonfresh state has \(W\le10^{-30}\), (T26)–(T28) give

\[
\mathcal Z\le10^{-16},\qquad
G/g\ge1-10^{-10},\qquad
K/g\le4\cdot10^{-6}.
\]

Therefore that state is accepted and its prescribed step is exactly \(1/2\).

Let \(n\) be the first index with \(W_n>10^{-30}\). The first \(n\) steps are consequently all half-steps. From (T15),

\[
W_n\le10^{-20}<1.
\]

Each such step increases \(\log W\) by at least \(1/4\), by (T32). Since
\(-\log W_0=3b+2\varepsilon\),

\[
n\le20b.
\tag{T36}
\]

The remaining part of the epoch has uniformly bounded logarithmic motion:

\[
\sum_{k=n}^{M-1}(2a_k+\beta_k)
=\log(W_{\rm end}/W_n)
<43\log10<150.
\tag{T37}
\]

Thus the tail contributes less than \(75\) to \(\sum a_k\) and less than \(150\) to \(\sum\beta_k\), regardless of its number or sizes of steps.

For the half-step prefix, the exact increments are

\[
a_k=\log(q_k+P)-\log(q_k+\delta),
\qquad
\beta_k=\log(q_k+\delta)-\log(q_k+S).
\tag{T38}
\]

The corresponding constants at \(q=1,P=\mathcal A_\sigma,S=0\) are precisely \(a_\sigma,\beta\) from (T5).

From (T9),

\[
|P-\mathcal A_\sigma|
\le2.2\cdot10^8e^{-b+100},
\qquad
S\le1.1\cdot10^8e^{-2b+200}.
\tag{T39}
\]

Also, throughout the half-step prefix,
\(\omega_{k+1}<.996\omega_k\), so (T13) gives

\[
\sum_{k<n}|q_k-1|
\le .011\sum_{k<n}\omega_k
\le2.75\omega_s
\le2.75\cdot10^{-6}.
\tag{T40}
\]

The relevant logarithmic derivatives in (T38) have absolute value at most \(2\). Equations (T36), (T39), and (T40), with \(b\ge1000\), therefore give

\[
\left|\sum_{k<n}a_k-na_\sigma\right|<10,
\qquad
\left|\sum_{k<n}\beta_k-n\beta\right|<10.
\tag{T41}
\]

These are explicit uniform bounds. The two residual-error sums are at most
\[
 E_P=20b(2.2\cdot10^8)e^{-b+100}<1,\qquad
 E_S=20b(1.1\cdot10^8)e^{-2b+200}<1.
\]
Both functions decrease for $b\ge1000$. At $b=1000$, replacing $e^{-900}$ and $e^{-1800}$ by the larger $2^{-900}$ and $2^{-1800}$ already gives the strict bounds. Each deviation $|a_k-a_\sigma|$ or $|\beta_k-\beta|$ from the corresponding limiting increment in (T38) is bounded by $4|q_k-1|$ plus twice its relevant residual error; hence each sum in (T41) is at most $4(2.75\cdot10^{-6})+2\max(E_P,E_S)<10$. This controls the entire half-step prefix uniformly in its length.

Put

\[
X=-\log(x/\omega),\qquad Y=-\log(y/\omega).
\]

Combining the prefix estimates with (T37) gives

\[
X_{\rm end}=2b+\varepsilon-na_\sigma+E_X,
\qquad
Y_{\rm end}=b+n\beta+E_Y,
\qquad
|E_X|,|E_Y|\le200.
\tag{T42}
\]

By (T17) and (T35),

\[
\varepsilon_{\rm end}:=Y_{\rm end}-2X_{\rm end}
=\log W_{\rm end},
\qquad |\varepsilon_{\rm end}|\le100.
\tag{T43}
\]

Eliminating \(n\) from (T42) yields

\[
\begin{aligned}
X_{\rm end}
={}&C_\sigma b
+\frac{\beta}{2a_\sigma+\beta}\varepsilon
-\frac{a_\sigma}{2a_\sigma+\beta}\varepsilon_{\rm end}\\
&+\frac{a_\sigma}{2a_\sigma+\beta}E_Y
+\frac{\beta}{2a_\sigma+\beta}E_X.
\end{aligned}
\tag{T44}
\]

Every displayed coefficient of an error has absolute value below \(1\). Thus

\[
|X_{\rm end}-C_\sigma b|\le1000.
\]

Together with (T43), this proves (T8) and completes the epoch lemma.



\subsection{The two-epoch map expands}

For the chosen constants,

\[
a_{\mathrm H}=\log\frac{1030401}{199},
\qquad
a_{\mathrm L}=\log\frac{2030301}{1990000},
\qquad
\beta=\log\frac{199}{100}.
\tag{T45}
\]

We need only the elementary bounds

\[
0<a_{\mathrm H}<9,\qquad
0<a_{\mathrm L}<.021,\qquad
\beta>.68.
\tag{T46}
\]

For the first, \(1030401/199<5200<e^9\).
For the second, use \(\log(1+x)<x\).
For the third, putting \(z=99/299\),

\[
\log\frac{199}{100}
=2\operatorname{arctanh}z
>2\left(z+\frac{z^3}{3}\right)
=\frac{18348264}{26730899}>.68.
\]

The function

\[
C(a,\beta)=\frac{a+2\beta}{2a+\beta}
\]

decreases in \(a>0\) and increases in \(\beta>0\). Therefore

\[
\begin{aligned}
C_{\mathrm H}C_{\mathrm L}
&>
\frac{9+1.36}{18+.68}
\frac{.021+1.36}{.042+.68}\\
&=\frac{357679}{337174}
>\frac{21}{20}=1.05.
\end{aligned}
\tag{T47}
\]

Also,

\[
\frac12<C_\sigma<2.
\tag{T48}
\]

Suppose a high-type epoch starts with parameter \(b\ge10^6\). The next, low-type epoch starts with parameter

\[
b'\ge C_{\mathrm H}b-1000
>\frac b2-1000\ge499000,
\tag{T49}
\]

so the epoch lemma remains applicable. The following high-type epoch starts with

\[
\begin{aligned}
b''
&\ge C_{\mathrm L}b'-1000\\
&\ge C_{\mathrm L}C_{\mathrm H}b
-1000C_{\mathrm L}-1000\\
&>1.05b-3000
\ge1.04b.
\end{aligned}
\tag{T50}
\]

Thus the high-type epoch parameters grow geometrically. Every intervening low-type epoch also remains within the lemma’s hypotheses. At each failed endpoint the actual gate provides a new error parameter with absolute value at most $100$; this parameter is reset by the endpoint identity and is not an error accumulated over all preceding epochs. Moreover $2b'+\varepsilon_{\rm next}>b'$ when $b'\ge1000$, so the two informative rows really exchange their small/large roles at every refresh.

This proves an infinite succession of finite epochs, alternating between the two informative rows.



\subsection{The uniform initialization enters this induction exactly}

For (T1),

\[
\frac{x_1}{\omega}=N^{-2},
\qquad
\frac{x_2}{\omega}=N^{-1}.
\]

The initial epoch is high-type and satisfies

\[
X_0=2\log N,\qquad Y_0=\log N.
\]

Thus its parameters are

\[
b_0=\log N=10^7\log2>5\cdot10^6,
\qquad \varepsilon_0=0.
\tag{T51}
\]

Moreover,

\[
\omega_0=\frac{N^2}{m}<\frac1N<10^{-6}.
\tag{T52}
\]

Equations (T51)–(T52) satisfy the hypotheses above. Consequently, the \textbf{exact uniformly initialized algorithm} has infinitely many refreshes, and the region

\[
x_1/\omega,\ x_2/\omega\le10^{-20}
\tag{T53}
\]

holds at every iterate.

All individual weights remain strictly positive: every update multiplier is at least \(1-t\ge1/2\). Every cache factorization is therefore well-defined.

Each epoch contains a fresh half-step, which decreases \(\omega\) by a factor below \(.996\); all other steps decrease it as well. Since there are infinitely many epochs,

\[
\omega_k\longrightarrow0.
\]

Equation (T53) then gives

\[
x_{1,k}\longrightarrow0,\qquad
x_{2,k}\longrightarrow0,\qquad
w_{0,k}\longrightarrow1.
\tag{T54}
\]

For every iterate, nonnegativity of the informative-pair contribution and evaluation at \(\theta=0\) give

\[
w_0\le v(w)
\le w_0+\delta\omega+B(x_1+x_2)
=1-.01\omega+.01(x_1+x_2)<1.
\tag{T55}
\]

Together with (T54),

\[
\boxed{\lim_{k\to\infty}v(w_k)=1<\frac{101}{100}=v^*.}
\tag{T56}
\]

The established gain bound, with \(g>0\), also shows that every accepted step strictly increases \(v\). Hence this is genuinely

\[
v(w_k)\uparrow1,
\]

not a failure of monotonicity or a finite-precision support deletion.

